\section{Packets from arbitrary orders}\label{sec:packets}

Fix a primitive totally real quartic field \(K\), an order \(\OO\subset K\),
and an ordering \(\sigma=(\sigma_1,\ldots,\sigma_4)\) of its real embeddings,
as in Theorem~\ref{thm:main}.  We describe the packet using finitely many
local lattice choices above each ideal class of the maximal order.  This
description lets us average a containment condition at each finite prime
before performing an integral at the real place.  The construction is
adapted from the packet model in
\cite[Section~2, Lemma~2.1]{PrimeDegreePackets}; we give the argument with
all sixteen sign patterns in degree four.

Throughout the paper put
\begin{equation}\label{eq:pkt-scale}
 R=\OO_K,\qquad D=|\Disc(K)|,\qquad q=[R:\OO],\qquad
 X=q\sqrt D,\qquad \kappa=\Res_{s=1}\zeta_K(s).
\end{equation}
In particular, \(|\Disc(\OO)|=q^2D=X^2\).  Absolute ideal norms are
extended multiplicatively to fractional \(R\)-ideals.  Thus
\begin{equation}\label{eq:pkt-covolume}
 \covol(\sigma(I))=\sqrt D\,\Nm I
\end{equation}
for every nonzero fractional \(R\)-ideal \(I\): this is the defining
discriminant identity for \(I=R\), follows by index for integral \(I\),
and then follows for fractional \(I\) after multiplication by an integer.
For a full \(\Z\)-lattice \(J\subset K\), write
\(J_p=J\otimes_\Z\Z_p\), and put
\[
 K_p=K\otimes_\Q\Q_p,\qquad R_p=R\otimes_\Z\Z_p,\qquad
 \OO_p=\OO\otimes_\Z\Z_p .
\]
Here \(K_p\) is a product of finite extensions of \(\Q_p\), and \(R_p\)
is the product of their integer rings.  We use the usual arithmetic of
fractional ideals in the Dedekind domain \(R\); see
\cite[Theorems~3.7, 3.20, and~3.29]{MilneANT}.

\subsection{The finite local choices}

Choose one fractional \(R\)-ideal \(I\) in each ordinary ideal class of
\(R\).  Write \(h_K\) for the number of these classes.  At every prime
choose a generator
\[
 I_p=\xi_p R_p,\qquad \xi_p\in K_p^\times.
\]
Such a generator exists because each component of \(R_p\) is a discrete
valuation ring.  Let \(\mathscr L_{\OO}(I)\) be the set of full lattices
\(L\subset I\) whose completions have the form
\begin{equation}\label{eq:pkt-local-choices}
 L_p=\xi_p t_p\OO_p,\qquad t_p\in R_p^\times
 \quad\text{for every }p.
\end{equation}
Replacing \(\xi_p\) by another generator merely changes \(t_p\) by a
unit, so this set does not depend on the chosen generators.

\begin{lemma}\label{lem:pkt-local-labels}
The set \(\mathscr L_{\OO}(I)\) is finite, and its cardinality is
\begin{equation}\label{eq:pkt-label-count}
 N_{\OO}=\prod_{p\mid q}[R_p^\times:\OO_p^\times],
\end{equation}
independently of \(I\).  Its members satisfy
\[
 [I:L]=q,\qquad RL=I,
\]
and are proper invertible fractional \(\OO\)-ideals.  Every ordinary
class of proper invertible fractional \(\OO\)-ideals is represented,
up to multiplication by \(K^\times\), by a member of one of the sets
\(\mathscr L_{\OO}(I)\).  Uniform probability on
\(\mathscr L_{\OO}(I)\) is induced by independent normalized Haar
probabilities for \(t_p\in R_p^\times\).
\end{lemma}

\begin{proof}
The additive group \(R/\OO\) has order \(q\), so \(qR\subset\OO\).
For every \(p\) and every unit \(t_p\), it follows that
\[
 qI_p=\xi_p t_p qR_p\subset \xi_p t_p\OO_p\subset I_p.
\]
At \(p\nmid q\) the middle lattice is \(I_p\).  At the remaining
primes its image is a submodule of the finite module \(I_p/qI_p\).
The primary decomposition of \(I/qI\) identifies its \(p\)-primary
part with \(I_p/qI_p\).  Thus a choice of these local submodules
patches to a unique subgroup of \(I/qI\), and its inverse image in
\(I\) is the unique lattice with the prescribed completions.

The stabilizer of \(\OO_p\) under multiplication by \(R_p^\times\)
is exactly \(\OO_p^\times\).  Indeed, \(t_p\OO_p=\OO_p\) implies
that both \(t_p\) and \(t_p^{-1}\) belong to \(\OO_p\), by applying
the equality and its inverse to \(1\).  The converse is immediate.
The local choices therefore form the finite homogeneous space
\(R_p^\times/\OO_p^\times\).  Its invariant probability is uniform
and is the image of normalized unit Haar measure.  Patching over the
finitely many primes dividing \(q\) proves
Equation~\eqref{eq:pkt-label-count} and the assertion about
probabilities.  The index of the lattice in
Equation~\eqref{eq:pkt-local-choices} is
\([R_p:\OO_p]\), and its \(R_p\)-span is \(I_p\) because \(1\in\OO_p\).
Taking products of the indices and comparing all completions gives
\([I:L]=q\) and \(RL=I\).

We record why the local condition describes exactly invertible ideals
of the order.  A full fractional \(\OO\)-ideal \(J\) is invertible
if and only if
\begin{equation}\label{eq:pkt-local-principality}
 J_p=\eta_p\OO_p\quad\text{for some }\eta_p\in K_p^\times
 \quad\text{at every }p.
\end{equation}
For the forward implication, complete an equality
\(JJ^{-1}=\OO\).  The resulting equality over \(\OO_p\) expresses
\(1\) as a finite sum \(\sum_i a_i b_i\), with \(a_i\in J_p\) and
\(b_i\in J_p^{-1}\).  At each maximal ideal \(\mathfrak m\) of the
semilocal ring \(\OO_p\), some \(a_i b_i\) is a unit after localization,
and that \(a_i\) generates \((J_p)_{\mathfrak m}\).  To obtain one
generator for all the maximal ideals, choose such an \(a_{\mathfrak m}\)
for each \(\mathfrak m\), and use the Chinese remainder theorem
\cite[Theorem~1.14]{MilneANT} to
choose \(c_{\mathfrak m}\in\OO_p\) congruent to \(1\) at
\(\mathfrak m\) and to \(0\) at every other maximal ideal.  Then
\(\eta_p=\sum_{\mathfrak m}c_{\mathfrak m}a_{\mathfrak m}\) generates
\((J_p)_{\mathfrak m}\) modulo \(\mathfrak m\) at each \(\mathfrak m\);
Nakayama's lemma \cite[Lemma~1.9]{MilneANT} shows that it generates each localization.  Hence it
generates \(J_p\).  Since \(J_p\) spans \(K_p\) over \(\Q_p\),
\(\eta_p\) has no zero component and lies in \(K_p^\times\).

Conversely, define
\[
 J^{-1}=\{a\in K:aJ\subset\OO\}.
\]
This definition commutes with completion: after choosing \(\Z\)-bases
of \(J\) and \(\OO\), the containment is a finite collection of
integrality conditions on the coordinates of multiplication by \(a\),
and the same conditions over \(\Z_p\) describe the completed lattice.
Under Equation~\eqref{eq:pkt-local-principality} we therefore have
\((J^{-1})_p=\eta_p^{-1}\OO_p\).  Consequently
\((JJ^{-1})_p=\OO_p\) for every \(p\), which implies \(JJ^{-1}=\OO\).
This proves the equivalence.  It also proves properness: if \(aJ\subset J\)
for an invertible \(J\), multiplication by \(J^{-1}\) gives
\(a\OO\subset\OO\), and hence \(a\in\OO\).  Thus its multiplier order
is \(\OO\).

The lattices in Equation~\eqref{eq:pkt-local-choices} are stable under
\(\OO_p\) at every prime, so comparison of completions shows that
\(\OO L\subset L\).  They are therefore fractional \(\OO\)-ideals
satisfying Equation~\eqref{eq:pkt-local-principality}.  In the other
direction, let \(J\) be an invertible fractional \(\OO\)-ideal.  Its \(R\)-span
\(RJ\) is a fractional \(R\)-ideal.  Multiplying \(J\) by an element
of \(K^\times\), we may arrange \(RJ=I\) for the chosen representative
of this maximal-order class.  If \(J_p=\eta_p\OO_p\), then
\(\eta_p R_p=I_p=\xi_pR_p\), so \(\eta_p=\xi_p t_p\) for a unit
\(t_p\in R_p^\times\).  Thus the resulting \(J\) belongs to
\(\mathscr L_{\OO}(I)\), as required.
\end{proof}

\subsection{Real parameters, identifications, and volume}

Let \(U=R^\times\), and let \(U^+\) be its totally positive subgroup.
Put
\[
 T_\infty=\{v=(v_1,\ldots,v_4)\in(\R^\times)^4:
                         \textstyle\prod_i|v_i|=1\}.
\]
We give \(T_\infty\) counting measure on its sixteen sign components
and measure \(du_1\,du_2\,du_3\) on the logarithms
\(u_i=\log|v_i|\), with \(u_4=-u_1-u_2-u_3\).  Denote this measure
by \(dv\).  The unit theorem
\cite[Theorems~5.1 and~5.9]{MilneANT} says that
\[
 \ell(U)=\{(\log|\sigma_i(\epsilon)|)_{i=1}^4:\epsilon\in U\}
\]
is a full lattice in the hyperplane \(\sum_i u_i=0\).  Its
covolume for the deleted-coordinate measure is the ordinary regulator
\(R_K\).  Its kernel on \(U\) consists of \(1\) and \(-1\), since
these are the roots of unity in a totally real field.

Let \(\mathscr F_U\) be a measurable fundamental domain for the embedded
group \(\sigma(U)\) in \(T_\infty\).  Reducing logarithms modulo
\(\ell(U)\) leaves sixteen sign choices, on which the remaining kernel
\(\{1,-1\}\) identifies pairs of opposite signs.  It follows that
\(\vol(\mathscr F_U)=8R_K\).  The analytic class-number formula in
signature \((4,0)\), with two roots of unity, gives
\begin{equation}\label{eq:pkt-unit-volume}
 h_K\vol(\mathscr F_U)=8h_KR_K=\sqrt D\,\kappa;
\end{equation}
see \cite[Chapter~V, Theorem~2.4(a)]{MilneCFT}.

For \(L\in\mathscr L_{\OO}(I)\), set
\[
 c_I=(X\Nm I)^{-1/4}.
\]
Equations~\eqref{eq:pkt-covolume} and \([I:L]=q\) show that
\(c_I\sigma(L)v\) has covolume one for every \(v\in T_\infty\).
This notation denotes the corresponding row lattice, and thus a point
of \(X_4\), after choosing a basis of determinant one.  This convention
also applies when the diagonal \(v\) has negative determinant.

\begin{proposition}[Packet model and volume]\label{prop:pkt-model}
The maps
\[
 (L,v)\longmapsto c_I\sigma(L)v
 \qquad
 (L\in\mathscr L_{\OO}(I),\ v\in T_\infty)
\]
give a bijection from the disjoint union over the chosen \(I\) of the
quotients by
\begin{equation}\label{eq:pkt-unit-action}
 \epsilon:(L,v)\longmapsto
       (\epsilon L,\sigma(\epsilon)^{-1}v),\qquad \epsilon\in U,
\end{equation}
onto \(\bigcup_{\mathcal Y\in\mathcal P_{\OO,\sigma}}\mathcal Y\subset X_4\).
Counting measure on the \(L\)'s times \(dv\) descends to the sum of
the Haar measures of those orbits.  In particular,
\begin{equation}\label{eq:pkt-total-volume}
 \sum_{\mathcal Y\in\mathcal P_{\OO,\sigma}}\vol(\mathcal Y)
       =N_{\OO}\sqrt D\,\kappa .
\end{equation}
Writing \(\mu_*=\mu_{\OO,\sigma}\), for every nonnegative Borel
function \(\Phi\) on \(X_4\) one has
\begin{equation}\label{eq:pkt-integration}
 \int_{X_4}\Phi\,d\mu_*
 =\frac{1}{h_KN_{\OO}\vol(\mathscr F_U)}
   \sum_I\sum_{L\in\mathscr L_{\OO}(I)}
       \int_{\mathscr F_U}\Phi(c_I\sigma(L)v)\,dv .
\end{equation}
The identity permits infinite values on both sides.
\end{proposition}

\begin{proof}
The action in Equation~\eqref{eq:pkt-unit-action} preserves the
parameter set, since \(\epsilon I=I\) for \(\epsilon\in U\), and
preserves its image lattice.  Lemma~\ref{lem:pkt-local-labels} shows
that every ordinary invertible \(\OO\)-ideal class occurs.  If two
ideals differ by multiplication by \(\alpha\in K^\times\), their
normalized embedded lattices differ by the diagonal with coordinates
\(\sigma_i(\alpha)/|\Nm_{K/\Q}(\alpha)|^{1/4}\), an element of
\(T_\infty\).  Finally \(T_\infty\) is the union of all sixteen
sign translates of \(A_4\).  These facts prove surjectivity.

Suppose two displayed normalized lattices coincide.  After absorbing
the ratio of their positive scalars, this equality gives a diagonal
map carrying \(\sigma(L)\) onto \(\sigma(L')\).  Choose \(0\ne b\in L\)
and let \(b'\in L'\) be its image.  None of the four coordinates of
\(\sigma(b)\) is zero.  The diagonal entries must therefore be
\(\sigma_i(b')/\sigma_i(b)=\sigma_i(b'/b)\).  The map is multiplication
by \(\alpha=b'/b\), and \(L'=\alpha L\).  Taking \(R\)-spans gives
\(I'=\alpha I\).  The chosen representatives then force \(I'=I\)
and \(\alpha I=I\), so \(\alpha\in U\).  The scalar normalizations
are now equal, and the two parameters differ by exactly
Equation~\eqref{eq:pkt-unit-action}.  This proves injectivity, including
all identifications involving signs.  Equality of covolume-one row
lattices is the same as equality in \(X_4\): any two positively oriented
bases of the same lattice differ by an element of \(\mathrm{SL}_4(\Z)\).

It remains to check that this identification gives the stipulated
orbit weights.  Write
\(\OO^{\times,+}=\OO^\times\cap U^+\).  The \(U\)-stabilizer of a
label \(L\) is \(\OO^\times\), by its multiplier order.  The stabilizer
of \(L\) together with a sign pattern is therefore
\(\OO^{\times,+}\).  This subgroup has finite index in \(U\).
For \(q=1\) this follows directly from \(\OO=R\).  For \(q>1\),
the kernel of \(U\longrightarrow(R/qR)^\times\) lies in
\(\OO^\times\), since both a unit congruent to \(1\) modulo \(qR\)
and its inverse belong to \(\OO\), and imposing positive signs costs
at most a further factor of sixteen.  The same diagonal-identification
argument shows that all periods of the corresponding \(A_4\)-orbit
are precisely \(\ell(\OO^{\times,+})\).  They form a full lattice,
so the orbit is compact.

Combining a label with a sign produces \(16N_{\OO}\) finite choices
above each \(I\).  Each \(U\)-orbit of these choices has
\([U:\OO^{\times,+}]\) elements and parametrizes one \(A_4\)-orbit.
On that orbit the quotient of logarithmic Lebesgue measure is its
Haar measure, with volume
\[
 \covol\bigl(\ell(\OO^{\times,+})\bigr)
 =[\ell(U):\ell(\OO^{\times,+})]R_K
 =\frac{[U:\OO^{\times,+}]}{2}R_K .
\]
The last equality uses \(\ker(\ell|_U)=\{1,-1\}\) and
\(-1\notin\OO^{\times,+}\).  Consequently the number of distinct
orbits times their common volume is
\[
 \frac{16h_KN_{\OO}}{[U:\OO^{\times,+}]}
 \cdot \frac{[U:\OO^{\times,+}]}2 R_K
 =8h_KN_{\OO}R_K.
\]
This proves both the assertion about the descended measure and
Equation~\eqref{eq:pkt-total-volume}, by
Equation~\eqref{eq:pkt-unit-volume}.

One can also take \(\mathscr L_{\OO}(I)\times\mathscr F_U\) as a
fundamental domain for the diagonal action in
Equation~\eqref{eq:pkt-unit-action}: reducing \(v\) into
\(\mathscr F_U\) transports \(L\) by the same unique unit, outside
the usual null boundary set.  Counting measure times \(dv\) is
invariant under this action.  These domains have total measure
\(h_KN_{\OO}\vol(\mathscr F_U)\), equal to the total volume just
computed.  Dividing their integral by that quantity proves
Equation~\eqref{eq:pkt-integration}.
\end{proof}
